\documentclass[11pt]{article}
\usepackage{amsmath,amssymb,amsfonts,geometry}
\usepackage{hyperref}
\geometry{margin=1in}

\title{Prime-Indexed Multiplicity Gravity:\\
Mother Theory, SUSY/SUGRA, and Dynamic Feedback}
\author{Citizen Gardens \\ \\ The Foundation of Multiplicity}
\date{April 2, 2026}

\begin{document}
\maketitle

\begin{abstract}
We develop a unified framework---\emph{prime-indexed multiplicity gravity}---that integrates Mother Theory's topological quantum gravity, supersymmetry (SUSY), supergravity (SUGRA), and Multiplicity Theory's eigenvalue dynamics into a single, recursively stable architecture. Mother Theory contributes topological invariants (Euler characteristic and Pontryagin index) and non-trivial boundary terms to a quantum-corrected Einstein sector.\,[file:1] Supersymmetry and supergravity supply the local field content and symmetries, organized as prime-indexed multiplicity-2 fibers (SUSY doublets) over spacetime.\,[file:1] Multiplicity Theory provides the spectral/dynamical engine via a global multiplicity functional \(M(t)\) and a delayed feedback law that can preserve or break SUSY fiber-by-fiber.\,[file:1] SUSY breaking appears microscopically as \emph{multiplicity splitting}: a prime-labeled eigenvalue of multiplicity two splits into two non-degenerate multiplicity-one modes when a prime-dependent load \(L_p(t)\) crosses a critical threshold.\,[file:1] Unbroken fibers correspond to stable dark-matter candidates, while broken fibers populate the visible sector. We give explicit \(2\times 2\) SUSY matrices, embed them into a 1+1D SUSY field theory with prime-indexed modes, and provide a cosmological evolution scenario for the loads \(L_p(t)\) in FRW and high-curvature backgrounds.\,[file:1]
\end{abstract}

\section{Executive Summary}

\subsection{Architecture}

The framework has four interacting layers:

\begin{itemize}
  \item \textbf{Geometric/topological layer (Mother Theory).} Spacetime dynamics are governed by quantum-corrected Einstein equations augmented with topological invariants: the Euler characteristic \(\chi(g)\) and Pontryagin index \(P(g)\) enter both the action and the field equations, refining the description of black holes, wormholes, and other non-trivial geometries.\,[file:1]
  \item \textbf{Symmetry/field-theory layer (SUSY/SUGRA).} Supersymmetry introduces supermultiplets of bosons and fermions, and supergravity promotes SUSY to a local symmetry including gravitons and gravitinos, yielding supersymmetric quantum gravity.\,[file:1]
  \item \textbf{Spectral/dynamical layer (Multiplicity).} A global multiplicity functional \(M(t)\) encodes interactions between eigenvalues \(\{\lambda_i\}\) via time-dependent weights \(w_i(t)\), phases \(\phi_i(t)\), and correlation factors.\,[file:1]
  \item \textbf{Feedback layer (fractal corrections and delays).} Delayed feedback equations adjust the weights \(w_i(t)\) based on past values of \(M(t)\), implementing scale-dependent, non-Markovian corrections interpreted as fractal or McGinty-type contributions to curvature and stress-energy.\,[file:1]
\end{itemize}

The novel step is to make the SUSY/SUGRA sector \emph{prime-indexed}: each SUSY doublet becomes a fiber \(\mathbb{C}_p^2\) labeled by a prime \(p\), with the global Hilbert space organized as a direct sum/tensor product of these fibers.\,[file:1] The same prime labels appear in the multiplicity functional and the feedback law, allowing SUSY breaking to be described as a controlled change in multiplicity.

\subsection{Prime-Indexed SUSY Fibers and Dark Matter}

At the microscopic level, each SUSY doublet is represented as a prime-indexed multiplicity-2 fiber with explicit \(2\times 2\) supercharge and Hamiltonian matrices obeying the SUSY algebra.\,[file:1] When the fiber is unbroken, its eigenvalue has multiplicity 2 and contributes coherently to \(M(t)\) as a single prime mode; when broken, the degeneracy is lifted and the eigenvalue splits into two multiplicity-1 contributions.\,[file:1]

A feedback-induced pitchfork bifurcation in the boson--fermion weight asymmetry \(\Delta_p(t)\) realizes SUSY breaking dynamically. Primes whose loads \(L_p(t)\) never cross the critical value \(L_p^{\mathrm{crit}}\) retain unbroken multiplicity-2 fibers and behave as stable dark-matter sectors, while primes that do cross \(L_p^{\mathrm{crit}}\) become visible-sector or intermediate dark states.\,[file:1]

\subsection{Field-Theoretic and Cosmological Realization}

The 0+1D SUSY fiber construction uplifts directly to a 1+1D SUSY field theory: each Fourier mode on a spatial circle is associated with a prime \(p_n\), and the corresponding mode operators form a prime-indexed SUSY doublet.\,[file:1] The multiplicity functional and feedback equations act on the mode weights \(w_{p_n\alpha}(t)\), allowing mode-by-mode SUSY breaking and multiplicity splitting.

Cosmologically, the prime loads \(L_p(t)\) depend on the FRW scale factor \(a(t)\) and the Supgra-corrected effective energy density stored in each mode. During inflation, many visible-sector primes cross their critical loads and split; dark-matter primes remain below threshold. After reheating, loads decay, freezing in a pattern of broken vs.\ unbroken primes that determines the partition between visible and dark sectors.\,[file:1] Near black-hole horizons and other high-curvature regions, local load spikes can trigger additional local multiplicity splitting for visible primes while leaving dark-matter primes intact.\,[file:1]

\section{Comprehensive Mathematical Overview}

\subsection{Mother Theory: Topological Quantum Gravity}

\subsubsection{Quantum-Corrected Einstein Equations with Topological Invariants}

Mother Theory augments the Einstein equations by incorporating quantum corrections and topological terms.\,[file:1] The quantum-corrected equations can be written schematically as
\begin{align}
R_{\mu\nu}
-\frac{1}{2} g_{\mu\nu} R
+\Lambda g_{\mu\nu}
+\Delta_{\mu\nu}\,\Omega_{\mu\nu}(u)\,Q_{\mu\nu}
+\epsilon\,\nabla^{2} Q_{\mu\nu}
&= 8\pi G \,\langle T_{\mu\nu}\rangle_{\mathrm{quantum}}
+\lambda_1 \chi(g)\,g_{\mu\nu}
+\lambda_2 P(g)\,g_{\mu\nu},
\end{align}
where:
\begin{itemize}
  \item \(R_{\mu\nu}\) is the Ricci tensor and \(R\) the Ricci scalar;\,[file:1]
  \item \(\Lambda\) is the cosmological constant;\,[file:1]
  \item \(\Delta_{\mu\nu}\Omega_{\mu\nu}(u)Q_{\mu\nu}\) and \(\epsilon\nabla^{2}Q_{\mu\nu}\) encode quantum corrections and fluctuations, significant near black holes and in the early universe;\,[file:1]
  \item \(\chi(g)\) is the Euler characteristic and \(P(g)\) the Pontryagin index of the spacetime manifold, treated as topological sources;\,[file:1]
  \item \(\langle T_{\mu\nu}\rangle_{\mathrm{quantum}}\) is the quantum-corrected stress-energy tensor.
\end{itemize}

\subsubsection{Action with Euler Characteristic and Pontryagin Index}

These topological contributions arise from an enhanced gravitational action of the form
\begin{align}
S_{\mathrm{Multiplicity}}
=
\int d^4x\,\sqrt{-g}\left[
\frac{1}{16\pi G}(R-2\Lambda)
+\lambda_1\,\chi(g)+\lambda_2\,P(g)
\right],
\end{align}
where \(\chi(g)\) encodes global connectivity and \(P(g)\) encodes global curvature and field connections.\,[file:1] Varying this action yields the topological terms in the field equations and allows the theory to model black holes, wormholes, cosmic strings, and other non-trivial spacetimes.\,[file:1]

\subsubsection{Boundary Terms and Higher-Dimensional Embedding}

Non-trivial boundary conditions are incorporated via the Gibbons--Hawking--York term\,[file:1]
\begin{align}
S_{\mathrm{boundary}}
=
\int_{\partial M} d^3x\,\sqrt{-h}\,\bigl(K+f(\Theta)\bigr),
\end{align}
where \(\partial M\) is the boundary, \(h\) is the induced metric, \(K\) the extrinsic curvature, and \(f(\Theta)\) encodes defects or twisting (e.g.\ cosmic strings or wormhole throats).\,[file:1] Mother Theory can also be embedded into higher-dimensional metrics of Kaluza--Klein type,\,[file:1]
\begin{align}
g_{AB}
=
\begin{pmatrix}
g_{\mu\nu}+\phi^{2} A_{\mu} A_{\nu} & \phi^{2} A_{\mu}\\
\phi^{2} A_{\nu} & \phi^{2}
\end{pmatrix},
\end{align}
where \(A_\mu\) is a vector potential and \(\phi\) a scalar controlling the extra dimension.

\subsection{Multiplicity Functional and Eigenvalue Dynamics}

\subsubsection{Global Multiplicity Functional}

Multiplicity Theory introduces a global functional \(M(t)\) over an eigenvalue spectrum \(\{\lambda_i\}\) with associated weights \(w_i(t)\), phases \(\phi_i(t)\), and correlation factors.\,[file:1] The general form is
\begin{align}
M(t)
&=
\sum_{i=1}^{N}\sum_{j=1}^{N}
M(\lambda_i,\lambda_j)\,
C_{ij}(t)\,
\gamma_{ij}(t)\,
\delta_{ij}(t)\,
w_i(t)\,
\cos\bigl(\phi_i(t)-\phi_j(t)\bigr),
\label{eq:Multiplicity}
\end{align}
where
\begin{itemize}
  \item \(M(\lambda_i,\lambda_j)\) encodes interactions between eigenvalues \(\lambda_i,\lambda_j\);\,[file:1]
  \item \(C_{ij}(t)\) is a correlation matrix capturing entanglement;\,[file:1]
  \item \(\gamma_{ij}(t)\) and \(\delta_{ij}(t)\) describe coherence and decoherence;\,[file:1]
  \item \(w_i(t)\) are dynamic weights, and \(\phi_i(t)\) phases associated with each eigenstate.\,[file:1]
\end{itemize}

\subsubsection{Prime-Indexed Reinterpretation}

We reinterpret each index \(i\) as a pair \((p,\alpha)\), where \(p\in\mathbb{P}\) is a prime and \(\alpha\in\{\mathrm{B},\mathrm{F}\}\) labels bosonic/fermionic components within a SUSY doublet.\,[file:1] Then \eqref{eq:Multiplicity} becomes
\begin{align}
M(t)
&=
\sum_{p,q\in\mathbb{P}}
\sum_{\alpha,\beta\in\{\mathrm{B},\mathrm{F}\}}
M(\lambda_{p\alpha},\lambda_{q\beta})\,
C_{(p\alpha)(q\beta)}(t)\,
\gamma_{(p\alpha)(q\beta)}(t)\,
\delta_{(p\alpha)(q\beta)}(t)\,
w_{p\alpha}(t)\,
\cos\bigl(\phi_{p\alpha}(t)-\phi_{q\beta}(t)\bigr).
\end{align}
This separates into intra-fiber terms (\(p=q\)) and inter-fiber terms (\(p\neq q\)), with each prime \(p\) labeling a SUSY fiber.

\subsubsection{Dynamic Feedback Law}

The evolution of the weights is governed by a delayed feedback equation\,[file:1]
\begin{align}
\frac{dw_i(t)}{dt}
=
F_w\bigl(M(t-\tau),M(t-\tau-\Delta t)\bigr),
\end{align}
which, in prime-indexed form, becomes a fiber-wise map
\begin{align}
\frac{d}{dt}
\begin{pmatrix}
w_{p\mathrm{B}}(t)\\
w_{p\mathrm{F}}(t)
\end{pmatrix}
=
\mathbf{F}_w^{(p)}\bigl(\mathbf{M}_p(t-\tau),\mathbf{M}_p(t-\tau-\Delta t)\bigr),
\end{align}
where \(\mathbf{M}_p\) summarizes the intra-fiber multiplicity contributions for prime \(p\).\,[file:1] This feedback can preserve or break SUSY on each fiber, depending on the load \(L_p\) and the form of \(\mathbf{F}_w^{(p)}\).

\subsection{SUSY as Prime-Indexed Multiplicity-2 Fibers}

\subsubsection{Microscopic \(2\times 2\) SUSY Matrices}

For each prime \(p\), define a local SUSY fiber \(\mathbb{C}_p^2\) with basis
\begin{align}
|p,\mathrm{B}\rangle =
\begin{pmatrix}1\\0\end{pmatrix},\qquad
|p,\mathrm{F}\rangle =
\begin{pmatrix}0\\1\end{pmatrix}.
\end{align}
On this fiber, the supercharges and Hamiltonian are
\begin{align}
Q_p
&=
\begin{pmatrix}
0 & 0\\
\sqrt{E_p} & 0
\end{pmatrix},
&
Q_p^\dagger
&=
\begin{pmatrix}
0 & \sqrt{E_p}\\
0 & 0
\end{pmatrix},
&
H_p
&=
E_p\,\mathbb{I}_2.
\end{align}
They satisfy the SUSY algebra
\begin{align}
\{Q_p,Q_p^\dagger\}
=
Q_p Q_p^\dagger + Q_p^\dagger Q_p
=
2H_p
=
2E_p\,\mathbb{I}_2,
\end{align}
so the eigenvalue \(E_p\) has \emph{multiplicity 2}.\,[file:1]

\subsubsection{Unbroken SUSY and Multiplicity-2 Contributions in \(M(t)\)}

In the unbroken SUSY regime for prime \(p\),
\begin{align}
\lambda_{p\mathrm{B}} = \lambda_{p\mathrm{F}} = E_p,\quad
w_{p\mathrm{B}}(t) = w_{p\mathrm{F}}(t) \equiv w_p(t),
\end{align}
with intra-fiber coherence
\begin{align}
C_{(p\mathrm{B})(p\mathrm{F})}(t)\approx 1,\quad
\phi_{p\mathrm{B}}(t)-\phi_{p\mathrm{F}}(t)\approx 0.
\end{align}
The contribution of this fiber to \(M(t)\) reduces to
\begin{align}
M_p(t)
\approx
2\,M(E_p,E_p)\,w_p(t),
\end{align}
so the SUSY doublet acts as a single prime-labeled mode with multiplicity factor \(2\).\,[file:1]

\subsubsection{SUSY Breaking as Multiplicity Splitting}

Introduce a symmetry-breaking perturbation to the fiber Hamiltonian:
\begin{align}
H_p'=
\begin{pmatrix}
E_p+\delta_p & 0\\
0 & E_p-\delta_p
\end{pmatrix},\qquad \delta_p\neq 0.
\end{align}
The eigenvalues split:
\begin{align}
\lambda_{p\mathrm{B}}=E_p+\delta_p,\quad \lambda_{p\mathrm{F}}=E_p-\delta_p,
\end{align}
and the weights and coherence decouple:
\begin{align}
w_{p\mathrm{B}}(t)\neq w_{p\mathrm{F}}(t),\quad
C_{(p\mathrm{B})(p\mathrm{F})}(t)\to 0,\quad
\Delta\phi_p(t)\equiv\phi_{p\mathrm{B}}(t)-\phi_{p\mathrm{F}}(t)\ \text{drifts}.
\end{align}
The intra-fiber multiplicity contribution becomes\,[file:1]
\begin{align}
M_p(t)
&=
M(E_p+\delta_p,E_p+\delta_p)\,w_{p\mathrm{B}}(t)
+
M(E_p-\delta_p,E_p-\delta_p)\,w_{p\mathrm{F}}(t)
\\
&\quad
+2\,M(E_p+\delta_p,E_p-\delta_p)\,
C_{(p\mathrm{B})(p\mathrm{F})}(t)\,
w_{p\mathrm{B}}(t)w_{p\mathrm{F}}(t)\,
\cos\bigl(\Delta\phi_p(t)\bigr).
\end{align}
In the fully broken limit, \(C_{(p\mathrm{B})(p\mathrm{F})}(t)\to 0\), and \(M_p(t)\) reduces to two independent multiplicity-1 contributions. Thus SUSY breaking appears microscopically as \emph{multiplicity splitting}: a prime-indexed multiplicity-2 eigenvalue \(E_p\) is replaced by two non-degenerate multiplicity-1 eigenvalues \(E_p\pm\delta_p\) with independent weights and phases.\,[file:1]

\subsection{Fiber-Wise Feedback and Bifurcation Mechanism}

\subsubsection{Asymmetry Variables}

Define total weight and asymmetry on each prime:
\begin{align}
W_p(t) &= w_{p\mathrm{B}}(t)+w_{p\mathrm{F}}(t),\\
\Delta_p(t) &= w_{p\mathrm{B}}(t)-w_{p\mathrm{F}}(t).
\end{align}
SUSY corresponds to \(\Delta_p=0\). A simple feedback core is
\begin{align}
\frac{dW_p}{dt} &= -\alpha_p\bigl(W_p-\bar{W}_p\bigr),\\
\frac{d\Delta_p}{dt} &= -\beta_p\,\Delta_p,
\end{align}
with \(\alpha_p,\beta_p>0\) and target total weight \(\bar{W}_p\).\,[file:1] This has a stable fixed point \(W_p=\bar{W}_p\), \(\Delta_p=0\), i.e.\ an unbroken SUSY doublet.

\subsubsection{Load-Dependent Pitchfork Bifurcation}

To realize SUSY breaking dynamically, introduce a load parameter
\begin{align}
L_p(t)\equiv \mathcal{N}_p\bigl(M_p(t-\tau),M_p(t-\tau-\Delta t)\bigr)
\end{align}
and modify the asymmetry evolution to
\begin{align}
\frac{d\Delta_p}{dt}
=
\bigl[\mu_p(L_p)-\beta_p\bigr]\Delta_p - \gamma_p\,\Delta_p^3,
\qquad \gamma_p>0,
\end{align}
with a linear load response
\begin{align}
\mu_p(L_p) = a_p L_p - b_p,\qquad a_p>0.
\end{align}
Then:
\begin{itemize}
  \item For \(L_p < L_p^{\mathrm{crit}}=(\beta_p+b_p)/a_p\), \(\mu_p(L_p)<\beta_p\), and \(\Delta_p=0\) is stable (unbroken SUSY).
  \item For \(L_p > L_p^{\mathrm{crit}}\), \(\mu_p(L_p)>\beta_p\), and \(\Delta_p=0\) becomes unstable, with two nonzero stable fixed points
  \begin{align}
  \Delta_p^\star = \pm \sqrt{\frac{\mu_p(L_p)-\beta_p}{\gamma_p}},
  \end{align}
  corresponding to SUSY breaking and multiplicity splitting.\,[file:1]
\end{itemize}
The fiber-wise feedback map \(\mathbf{F}_w^{(p)}\) thus implements a pitchfork bifurcation: below critical load the SUSY doublet remains intact; above it, bosonic and fermionic weights diverge.

\subsection{Dark Matter as Unbroken Prime-Indexed SUSY Fibers}

\subsubsection{Stable Multiplicity-2 Dark Sector}

Dark matter arises naturally from primes \(p_{\mathrm{DM}}\) whose load never exceeds the critical value:
\begin{align}
L_{p_{\mathrm{DM}}}(t) < L_{p_{\mathrm{DM}}}^{\mathrm{crit}}\qquad \forall t.
\end{align}
For these primes, \(\Delta_{p_{\mathrm{DM}}}(t)=0\) remains stable, and the corresponding SUSY fibers stay unbroken multiplicity-2 doublets.\,[file:1] The weights satisfy \(w_{p_{\mathrm{DM}}\mathrm{B}}(t)=w_{p_{\mathrm{DM}}\mathrm{F}}(t)\equiv w_{p_{\mathrm{DM}}}(t)\), and the fiber contributes coherently to \(M(t)\) as
\begin{align}
M_{p_{\mathrm{DM}}}(t)\approx 2\,M(E_{p_{\mathrm{DM}}},E_{p_{\mathrm{DM}}})\,w_{p_{\mathrm{DM}}}(t).
\end{align}
These unbroken fibers correspond to stable, weakly coupled SUSY particles (neutralino- or gravitino-like) that gravitate but do not radiate in visible channels.\,[file:1]

\subsubsection{Supgra Action and Topological Counts}

The dark-matter sector appears in the Supgra-corrected action as an additional sum over unbroken fibers:
\begin{align}
S_{\mathrm{DM}}
=
\sum_{p_{\mathrm{DM}}}
\int d^4x\,\sqrt{-g}\;
\mathcal{L}_{\mathrm{SUSY},p_{\mathrm{DM}}},
\end{align}
where each \(\mathcal{L}_{\mathrm{SUSY},p_{\mathrm{DM}}}\) is a standard SUSY/SUGRA Lagrangian on the locked fiber.\,[file:1] Topological invariants count these multiplicity-2 contributions via
\begin{align}
\chi(g)\supset\sum_{p_{\mathrm{DM}}}(-1)^{s(p_{\mathrm{DM}})}\cdot 2,
\end{align}
reflecting their spectral multiplicity.\,[file:1]

\subsection{1+1D SUSY Field-Theory Realization}

\subsubsection{Prime-Indexed Mode Expansion}

Consider a 1+1D SUSY scalar--fermion system on a spatial circle of length \(L\).\,[file:1] Expand the fields as
\begin{align}
\phi(x,t) &= \sum_{n\in\mathbb{Z}}\phi_n(t)\,e^{ik_n x},
&
\psi(x,t) &= \sum_{n\in\mathbb{Z}}\psi_n(t)\,e^{ik_n x},
\end{align}
with \(k_n = 2\pi n/L\). Associate each mode index \(n\) with a prime \(p_n\), forming a prime-indexed SUSY fiber \(\mathbb{C}_{p_n}^2\) with
\begin{align}
|p_n,\mathrm{B}\rangle \leftrightarrow \phi_n,\qquad
|p_n,\mathrm{F}\rangle \leftrightarrow \psi_n.
\end{align}
The fiber operators are as above, with mode energy \(E_{p_n} = \sqrt{k_n^2+m^2}\).\,[file:1] The full Hilbert space is
\begin{align}
\mathcal{H} = \bigotimes_n \mathbb{C}_{p_n}^2.
\end{align}

\subsubsection{Field-Theoretic Multiplicity Functional}

The multiplicity functional extends to field modes as
\begin{align}
M(t)
&=
\sum_{n,m}
\sum_{\alpha,\beta\in\{\mathrm{B},\mathrm{F}\}}
M(\lambda_{p_n\alpha},\lambda_{p_m\beta})\,
C_{(p_n\alpha)(p_m\beta)}(t)\,
\gamma_{(p_n\alpha)(p_m\beta)}(t)\,
\delta_{(p_n\alpha)(p_m\beta)}(t)\,
w_{p_n\alpha}(t)\,
\cos\bigl(\phi_{p_n\alpha}(t)-\phi_{p_m\beta}(t)\bigr),
\end{align}
with intra-mode terms (\(n=m\)) mirroring the 0+1D fibers and inter-mode terms (\(n\neq m\)) encoding momentum/prime couplings.\,[file:1] The fiber-wise feedback law acts on each mode:
\begin{align}
\frac{d}{dt}
\begin{pmatrix}
w_{p_n\mathrm{B}}(t)\\
w_{p_n\mathrm{F}}(t)
\end{pmatrix}
=
\mathbf{F}_w^{(p_n)}\bigl(\mathbf{M}_{p_n}(t-\tau),\mathbf{M}_{p_n}(t-\tau-\Delta t)\bigr).
\end{align}

\subsubsection{Coupling to Supgra and Mother Theory}

The Supgra-corrected Lagrangian decomposes as
\begin{align}
\mathcal{L}_{\mathrm{Supgra}}
=
\sum_n
\mathcal{L}_{\mathrm{SUSY},p_n}(\phi_n,\psi_n;g_{\mu\nu},\psi_\mu,\dots),
\end{align}
with masses and couplings allowed to depend on \(w_{p_n\alpha}(t)\) and fractal corrections.\,[file:1] Topological invariants can be expressed as sums over primes and multiplicities:
\begin{align}
\chi(g) = \sum_n (-1)^{s(p_n)}\,m_{p_n},\qquad
m_{p_n} =
\begin{cases}
2,& \text{unbroken SUSY mode,}\\
1+1,& \text{broken SUSY mode,}
\end{cases}
\end{align}
while the Pontryagin index \(P(g)\) can be related to phases and correlations across modes.\,[file:1]

\subsection{Cosmological Evolution of Prime-Indexed Loads}

\subsubsection{FRW Background and Load Scaling}

In an FRW spacetime with metric
\begin{align}
ds^2 = -dt^2 + a^2(t)\,d\vec{x}^{\,2},
\end{align}
the load controlling bifurcation on each prime is modeled as
\begin{align}
L_p(t) \equiv \mathcal{N}_p\bigl(M_p(t-\tau),M_p(t-\tau-\Delta t)\bigr)
\propto a^{-\sigma_p}(t)\,\rho_{\mathrm{eff}}(t;k_p),
\end{align}
where \(\rho_{\mathrm{eff}}(t;k_p)\) is the effective energy density in the mode with comoving wavenumber \(k_p\) and \(\sigma_p\) encodes redshift behavior.\,[file:1]

\subsubsection{Inflation, Reheating, and Selection of Primes}

During inflation, \(a(t)\sim e^{H_{\mathrm{inf}} t}\), and modes with \(k_p/a(t)\gg H_{\mathrm{inf}}\) are highly excited, leading to large loads \(L_p(t)\) and early SUSY breaking for many visible-sector primes.\,[file:1] Modes with \(k_p/a(t)\ll H_{\mathrm{inf}}\) freeze out with nearly constant loads.

After reheating, radiation and matter eras cause loads to decay,
\begin{align}
L_p(t)\to 0 \quad \text{as } t\to\infty,
\end{align}
freezing in a pattern where primes that crossed \(L_p^{\mathrm{crit}}\) are broken, while primes that never did remain unbroken dark-matter fibers.\,[file:1]

\subsubsection{High-Curvature Regions and Local Splitting}

In strongly curved regions, such as near black-hole horizons, the quantum-corrected Einstein equations with topological and fractal terms generate large local fluctuations in \(\langle T_{\mu\nu}\rangle_{\mathrm{Supgra}}\) and hence local spikes in \(L_p(t,\vec{x})\).\,[file:1] Visible primes can transiently exceed \(L_p^{\mathrm{crit}}\), triggering local SUSY breaking and multiplicity splitting, with implications for Hawking radiation and black-hole evaporation spectra.\,[file:1] Dark-matter primes are protected by parameter choices in \(\mu_p(L_p)=a_p L_p-b_p\) so that even extreme curvature does not push them past the bifurcation threshold.\,[file:1]

\section{Conclusion}

We have constructed a prime-indexed multiplicity framework that unifies Mother Theory's topological quantum gravity, SUSY/SUGRA field theory, and Multiplicity Theory's spectral dynamics into a single, recursively adaptive architecture.\,[file:1] At the microscopic level, SUSY doublets are realized as prime-indexed multiplicity-2 fibers with explicit \(2\times 2\) supercharge and Hamiltonian matrices; SUSY breaking emerges as multiplicity splitting when a load-dependent pitchfork bifurcation destabilizes the symmetric fiber configuration.\,[file:1] At the field-theoretic level, a 1+1D SUSY field theory can be organized as a direct sum of prime-indexed modes, with the multiplicity functional \(M(t)\) and feedback map acting mode by mode.\,[file:1]

Cosmologically, FRW expansion and high-curvature regions determine the time evolution of prime loads \(L_p(t)\), selecting which primes split (visible sector) and which remain unbroken (dark matter).\,[file:1] This provides a coherent path from microscopic multiplicity dynamics to macroscopic phenomena such as black-hole evaporation, inflation, and dark matter. The framework is amenable to numerical experiments with finite sets of primes and modes, offering a concrete route to testing its qualitative predictions.

\end{document}